% Chapter 6 — State: The Three Homes (KLEPPMANN). Budget 5 pp. Discharges §7.
\section{State: The Three Homes}\label{ch:homes}

This chapter discharges \S 7 of the constitution.

A unit's state is not one thing. Some of it is a property of the unit alone --- the
settlement level a future has printed, the fixings a variance swap has accrued, whether a
note's barrier has knocked --- the same fact for every holder. Some of it is a joint
property of the unit and a particular wallet --- the variation margin this holder applied
yesterday, the dividend that holder is owed. And some of it is the contractual text
itself --- a multiplier, a strike, a fixing schedule --- versioned and append-only. Three
kinds of fact, and the state model gives each its own home: \textbf{UnitStatus},
\textbf{PositionState}, and \textbf{ProductTerms}. With the balances of Chapter~3
these are the four components of the ledger's state, and the four kinds of change a
transaction bundles (Chapter~3) write them one for one.

\subsection{The three homes}

\textbf{ProductTerms} holds the versioned, append-only contractual terms of an
instrument: the future's multiplier, the option's barrier and payout, the variance swap's
strike and 252-fixing schedule. A new version is appended, never an edit; the current
terms are the latest version, and any past state is read against the version in force
then. Its facts are properties of the instrument's terms, and a corporate action that
changes those terms writes a new version (the split below).

\textbf{UnitStatus} holds the facts that are properties of the unit as a whole --- the
same for every holder: the last recorded settlement level, the running accrued sum of
squared returns, the live / triggered / expired lifecycle node. It is a materialised
projection of the unit-state changes on the log, written by the instrument's contract as
its life unfolds.

\textbf{PositionState} holds the facts that are joint properties of a unit and a wallet:
this holder's last applied settlement level, this holder's recorded dividend entitlement.
A fact lands here exactly when it depends on which wallet holds the unit --- when one
unit-level event yields a different value to each holder.

None of the three holds a balance --- that is the coordinate vector of
Chapter~3, the transaction's fourth kind of change --- and none holds a
view. The keys distinguish them: ProductTerms and UnitStatus are keyed by the unit,
PositionState by the wallet-and-unit pair.

\begin{lstlisting}
data Homes = Homes
  { productTerms  :: Unit           -> [TermsVersion]  -- append-only; latest is current
  , unitStatus    :: Unit           -> UnitFacts       -- properties of the unit alone
  , positionState :: (Wallet, Unit) -> PosFacts        -- joint unit-and-wallet facts
  }
-- each field is a fold of the log; discard and replay reconstructs it.
\end{lstlisting}

Two disciplines make the homes trustworthy, and both are checkable.

\begin{principle}[One fact, one home]\label{prin:one-home}
Every non-balance fact has exactly one home, determined by what the fact is a property of:
the unit alone (UnitStatus), the unit-and-wallet pair (PositionState), or the contract's
terms (ProductTerms). No fact is kept in two homes, so no two homes can disagree about it.
\end{principle}

\begin{invariant}[One legal writer]\label{inv:one-writer}
Every non-balance fact has exactly one legal writer --- the smart contract of the
instrument the fact belongs to --- and no other contract may write it. The Transaction
Executor enforces this at the door: a transaction that writes a fact outside its writer's
authority is refused. The invariant joins the structural catalogue of
Chapter~14 with constitution~\S 11's writer discipline.
\end{invariant}

\subsection{One fact, one home: the future exercises all three}

The daily life of a cash-settled future carries all three kinds of fact, so it visits all
three homes. FUT-IDX is a future on index IDX, multiplier~10, USD, settled-to-market;
W-ALPHA holds one contract. On day~1 the official close prints $1{,}000.00$, the Event
Monitor emits, and the future's contract fires. Three facts are in play, and each has one
home.

\begin{center}
\begin{tabular}{llll}
\toprule
Fact & Property of & Home & Legal writer \\
\midrule
Settlement print $1{,}000.00$        & the unit alone         & UnitStatus    & the print firing \\
W-ALPHA's applied level; VM $\USD{50}$ & the unit--wallet pair & PositionState & the print firing \\
Multiplier $10$                      & the contract's terms   & ProductTerms  & registration (v.\,1) \\
\bottomrule
\end{tabular}
\end{center}

The print is one number for every holder: it is a UnitStatus fact. The variation margin
depends on the level W-ALPHA last applied, a fact of the wallet-and-unit pair: it is a
PositionState fact, and the same print yields a different margin to a holder who entered
at a different level. The multiplier was fixed at registration in ProductTerms and changes
only by a new version. The print firing writes the first two facts in one transaction ---
two facts, two homes, one writer --- and reads yesterday's applied level back from
PositionState rather than holding it (Chapter~5).

\emph{Each fact has exactly one home, fixed by what it is a property of; the future,
carrying three kinds of fact, exercises all three homes.}

\subsection{The homes suffice}

The homes are not bookkeeping conveniences. Their reason for existing is a sufficiency
requirement, and it is the load the rest of the architecture rests on.

\begin{principle}[The homes suffice]\label{prin:sufficiency}
The three homes, together with the balances, carry all the information the smart contracts
require in order to execute. Every fact a firing must carry forward to a later firing is
written into a home by an earlier firing, so the later firing finds its entire context on
the record; if a contract would ever need to remember something, that something has a
home.
\end{principle}

This is the precise reason a contract is stateless and never remembers
(Chapter~5): statelessness is affordable only because every fact a firing
would want to remember already has a home the firing can read. The sufficiency is what
turns ``never remembers'' from an aspiration into a consequence.

Minimality is its other half. A home holds a \emph{fact} --- an input a later firing reads
--- never a \emph{view}. NAV, profit and loss, availability, and encumbrance are
projections, computed on demand from the homes and the log and never materialised into one
(Chapters~7 and~12). The homes carry every fact a contract
must read and no figure a valuation or a report recomputes --- which is exactly why no
store of ``realised PnL'' exists anywhere to be reconciled against the position record.

The variance swap exercises the requirement at full stretch.

\medskip\noindent\textbf{Episode --- the variance swap (T5): sufficiency, exercised.}
\emph{Trigger.} The 127th scheduled IDX close is recorded and the Event Monitor emits
VS-1's 127th fixing. VS-1 is a one-year OTC variance swap on IDX, 252 scheduled fixings
against the recorded IDX official closes --- the same series FUT-IDX settles on.
\emph{Contract.} VS-1's contract fires. To extend the realised variance it needs the whole
path so far: the previous close and the running sum of squared returns over fixings
$1,\dots,126$. It holds neither, and it is forbidden to.
\emph{What it finds on the record.} The previous close is a recorded observation. The
accrued statistic is a UnitStatus fact of VS-1 --- a property of the unit alone, the
realised path as the swap has seen it --- written by the 126th firing for exactly this
reader. At the mid-life checkpoint (fixing~126, half of the 252 fixings elapsed) it stands
at $A_{126}=1{,}805{,}000$, an integer at scale $10^{-8}$: the 126 accrued squared returns
compressed into one number. The firing reads $A_{126}$ and the previous close, forms the
127th return, adds its square, and writes $A_{127}$ into UnitStatus for the 128th firing.
\emph{The coupling.} The IDX close at fixing~126 is $1{,}005.00$ --- the same recorded
close FUT-IDX takes as its day-3 final settlement; fixings $124,125,126$ are the future's
three settlement days, closes $1{,}000$, $990$, $1{,}005$. One recorded observation, two
consumers: the future writes it as a settlement level into FUT-IDX's UnitStatus, the swap
folds it into VS-1's accrued statistic in VS-1's UnitStatus, and neither reads the other's
home.
\emph{Design point.} \emph{Firing~127 finds its entire context on the record --- 126
fixings compressed into one UnitStatus integer $A_{126}=1{,}805{,}000$ --- so a contract
that spans a year of history remembers nothing between firings. The homes suffice.}

\subsection{Two more facts finding their homes}

The dividend and the split add one home each: the entitlement is per-holder, the split is a
terms change.

\medskip\noindent\textbf{Episode --- the dividend (T3): an entitlement in PositionState.}
On 2026-05-15, the record date, ACME's contract fires and reads the owned positions as the
ledger then stands: W-ALPHA holds $10{,}000$ ACME. It records the entitlement
$10{,}000\times\USD{1.50}=\USD{15{,}000.00}$. A dividend entitlement is a property of the
unit-and-wallet pair --- the same $\USD{1.50}$ per share yields a different amount to each
holder --- so it lands in PositionState. The transaction is moveless: no cash moves on the
record date, but the position state now carries what W-ALPHA is owed, for the payment-date
firing to discharge (Chapter~5).
\emph{Design point.} \emph{An entitlement owed but unpaid is a fact with a home; the record
date writes it into PositionState movelessly, and the payment date reads it back.}

\medskip\noindent\textbf{Episode --- the split (T4): a new version in ProductTerms.}
On 2026-06-01 ACME's 2-for-1 split becomes effective. In one atomic transaction the holding
moves from $10{,}000$ to $20{,}000$ shares (a paired-leg move, Chapter~3) and
ACME's terms gain a new version: ProductTerms ACME version~2 carries the split, and the
recorded price frame reads $100.00$ as $50.00$ after it (the market data operator's work,
Chapter~8). A split is a property of the contract's terms, so
it lands in ProductTerms and nowhere else. The terms are not edited; a version is appended.
\emph{Design point.} \emph{Terms change only by appending a version; the split is one
ProductTerms fact, atomic with the move it accompanies, and the append-only discipline is
what lets any past state be read against the terms in force then.}

\subsection{Rebuildable by replay}

None of the three homes is an independent store. Each is a deterministic projection of the
immutable log, so it can be discarded and rebuilt at any time without loss.

\begin{principle}[Homes are projections]\label{prin:rebuildable}
Each home is a deterministic projection of the log --- ProductTerms of the terms changes,
UnitStatus of the unit-state changes, PositionState of the position-state changes. Any
home may be discarded and rebuilt by replaying the log, without loss of information, so no
home is a store that can drift from the record.
\end{principle}

This is the precise sense in which the homes are \emph{materialised projections} and not
second copies. A second copy is dangerous because two stores of one fact, maintained by
independent logic, eventually disagree --- the failure the framework exists to remove
(Chapter~1). A home cannot play that role: it has exactly one legal
writer (the one-writer invariant above), it holds each fact in exactly one place
(one fact, one home), and it is a pure fold of the one log everything else also
folds. Materialising the accrued statistic into UnitStatus is therefore not the
reintroduction of a store to reconcile; it is a cached read of the log, provably equal to
the log it came from. What is materialised is a fact a contract must read; what is never
materialised is a view a projection can recompute. Between them the homes carry everything
the contracts need and nothing they do not.
